% Generated draft: requires mathematical and source review.
\documentclass[11pt]{article}
\usepackage[margin=25mm]{geometry}
\usepackage{fontspec}
\setmainfont{Latin Modern Roman}
\usepackage{amsmath,amssymb,mathtools,unicode-math}
\setmathfont{Latin Modern Math}
\usepackage{enumitem,needspace,xurl,hyperref}
\setlength{\parindent}{0pt}
\setlength{\parskip}{5pt}
\setlength{\emergencystretch}{3em}
\newcommand{\N}{\mathbb{N}}
\newcommand{\Z}{\mathbb{Z}}
\newcommand{\Q}{\mathbb{Q}}
\newcommand{\R}{\mathbb{R}}
\newcommand{\C}{\mathbb{C}}
\newcommand{\F}{\mathbb{F}}
\newcommand{\A}{\mathbb{A}}
\newcommand{\PP}{\mathbb{P}}
\newcommand{\E}{\mathbb{E}}
\newcommand{\eps}{\varepsilon}
\newcommand{\dd}{\,\mathrm{d}}
\begin{document}
\section*{UnsolvedMath ID: 20002906}
\textbf{Source problem code:} catalogue entry
\subsection{Definitions and mathematical statement}
For integers \(0\le a\le b\le n\), let
\[
\operatorname{Fl}(a,b;n)=\{(A,B):A\subseteq B\subseteq\C^n,\ \dim A=a,\ \dim B=b\}.
\]
Its integral cohomology Schubert basis is indexed by words \(u\) containing \(a\) zeros, \(b-a\) ones and \(n-b\) twos, using a fixed Schubert convention. Define the structure constants by
\[
\sigma_u\sigma_v=\sum_w c_{u,v}^{w}\sigma_w.
\]
The supplied aid interprets the AIM request for Buch's weakened version of Knutson's conjecture as the two-step puzzle rule: \(c_{u,v}^{w}\) is the number of two-step puzzles with the boundary data corresponding to \(u,v,w\), under the rule's prescribed pieces and orientations. This is a statement about all two-step flag varieties, rather than the false unrestricted full-flag conjecture. The record does not list the puzzle pieces or specify boundary reversal/duality conventions; consequently it does not by itself define a self-contained counting function. An exact puzzle formulation must import those conventions from the intended theorem, not silently choose them.

The research aid reports the two-step rule proved by Buch, Kresch, Purbhoo and Tamvakis, and an equivariant version proved by Buch. These are reported literature results. The ordinary and equivariant counts should not be conflated: the latter require their equivariant weights.

A further reported support check uses the forgetful maps \(\pi_a:\operatorname{Fl}(a,b;n)\to\operatorname{Gr}(a,n)\) and \(\pi_b:\operatorname{Fl}(a,b;n)\to\operatorname{Gr}(b,n)\). Index a Grassmannian basis by \(02\)-words with respectively \(a\) or \(b\) zeros. In the standard stable-order lift convention:
\begin{itemize}
\item \(\ell_a\) keeps every zero and replaces the first \(b-a\) twos, scanning left to right, by ones.
\item \(\ell_b\) keeps every two and replaces the zeros after the first \(a\) zeros, scanning left to right, by ones.
\end{itemize}
With compatible Schubert conventions, \(\pi_j^*\sigma_\alpha=\sigma_{\ell_j(\alpha)}\) for \(j=a,b\). If \(c_{\alpha,\beta}^{\gamma}(\operatorname{Gr}(j,n))\) are ordinary Grassmannian Littlewood--Richardson coefficients, then
\[
c_{\ell_j(\alpha),\ell_j(\beta)}^{w}=
\begin{cases}
c_{\alpha,\beta}^{\gamma}(\operatorname{Gr}(j,n)),&w=\ell_j(\gamma),\\
0,&w\notin\operatorname{im}\ell_j.
\end{cases}
\]
This check applies when both inputs come from the same forgetful projection. It is not a rule for arbitrary pairs of two-step classes.

The source also motivates the two-step case by a quantum-to-classical reduction for Grassmannians. It prints a Buch--Kresch--Tamvakis--Yong citation in Duke Mathematical Journal 122 (2004), 125--143. The research aid explicitly says this citation is incorrect for that reduction and identifies the intended reduction as the 2003 JAMS kernel-and-span theorem of Buch, Kresch and Tamvakis. Both the original citation and the reported correction are preserved below; neither is silently substituted for a fully verified bibliography.
\subsection{Short English statement}
Express two-step flag Schubert constants by the intended puzzle rule, with the source's full-flag warning retained. The supplied aid reports that rule solved and gives exact Grassmannian-lift support checks, while also flagging the source's quantum-to-classical citation as incorrect.
\subsection{Sources}
\begin{itemize}
\item AIM workshop \emph{Generalized Kostka polynomials}, Problem (11), \url{https://aimath.org/WWN/kostka/prob2.pdf}.
\item Original printed reference: Buch, Kresch, Tamvakis and Yong, Duke Mathematical Journal 122 (2004), 125--143. The aid says this is not the stated quantum-to-classical reduction; its proposed correction is Buch, Kresch and Tamvakis, kernel-and-span theorem, JAMS (2003), without a complete title or page range supplied here.
\item The research summary attributes the two-step puzzle theorem to Buch, Kresch, Purbhoo and Tamvakis, and the equivariant theorem to Buch. Full bibliographic entries and the actual puzzle pieces are not supplied.
\item Local UnsolvedMath registry, permanent ID 20002906, exact code \texttt{catalogue entry}; canonical location \texttt{aim-special-functions-notes.json}, notes[10], dated 14 August 2026; statement labeled reconstructed and unverified, registry status solved and research classification solved in literature.
\item Research attempt 1: partial result, problem solved at the run, rigor label proved, novelty confidence low. The literature identification, citation correction and lift lemma are machine-generated research aids, not peer reviewed; independent verification is required before citation.
\end{itemize}
\end{document}
